\documentclass[11pt,a4paper]{article}
\usepackage{curso}
\begin{document}
\portada{Ejercicios --- Lección}{3.5}{Diseño de protocolos y aceleración de reloj}
\section{Ejercicios}
\begin{enumerate}[label=\textbf{E\arabic*.}]
\item (B) (Pensum C.4) ¿Qué cantidades cambian y cuáles no con la aceleración de reloj?
\item (B) Demuestre el cambio de variable $s=\lambda t$ en la ecuación escalada.
\item (B) Trabajo de una carga súbita $\kappa:1\to4$ desde equilibrio ($\kB T=1$).
\item (B) Con $\lambda=3$: ¿calor, potencia, $\Sigma$, $\eta$?
\item (I) Muestre que escalar solo $V\to\lambda V$ sin cambiar el ruido produce la distribución estacionaria a temperatura $T/\lambda$.
\item (I) Simule la rampa $\kappa:1\to4$ y ajuste $a$ en $W_{\text{dis}}\approx a/\tau$ para $\tau\ge2$.
\item (A) Para la rampa lineal de una trampa armónica, derive la ecuación de la varianza $\dot v=-2\kappa(t)v+2$ y calcule $\langle W\rangle(\tau)$ numéricamente.
\end{enumerate}
\section{Soluciones}
\begin{solucion}
\textbf{E1.} Calor y energías $\times\lambda$, duración $\div\lambda$, potencia $\times\lambda^2$; $\Sigma$, $\eta$ y las trayectorias (en tiempo reescalado) no cambian.

\textbf{E2.} $dt=ds/\lambda$ y $dW(t)\overset{d}{=}\lambda^{-1/2}d\tilde W(s)$: $-\mu\lambda V'ds/\lambda+\sqrt{2\mu\lambda\kB T}\lambda^{-1/2}d\tilde W$.

\textbf{E3.} $\frac12(4-1)\cdot1=1{,}5$.

\textbf{E4.} Calor ×3, potencia ×9, $\Sigma$ y $\eta$ iguales.

\textbf{E5.} $p\propto e^{-\lambda V/\kB T}=e^{-V/(\kB T/\lambda)}$.

\textbf{E6.} Con los datos de la lección, $W_{\text{dis}}(2)=0{,}14$, $(4)=0{,}077$, $(8)=0{,}041$: $a\approx0{,}3$.

\textbf{E7.} $d\langle x^2\rangle/dt=-2\kappa\langle x^2\rangle+2$ (Itô, $\mu=\kB T=1$) y $\langle W\rangle=\int_0^\tau\frac12\dot\kappa\,v\,dt$; integrando numéricamente se reproducen los valores de la tabla.
\end{solucion}
\end{document}
